%!TEX root=AImath.v4.tex

\chapter{ 人工智能的数学工具}
 \begin{introduction}
  \item Definition of Theorem
  \item Ask for help
  \item Optimization Problem
  \item Property of Cauchy Series
  \item Angle of Corner
\end{introduction}


微积分

- 导数和偏导数：用于计算函数的变化率，广泛应用于梯度下降算法和神经网络的训练。
- 链式法则：用于反向传播算法，以计算神经网络中的梯度。
- 梯度和Hessian矩阵：用于优化问题，特别是在高维空间中。
- 积分和多重积分：用于概率计算和期望值计算，如在贝叶斯推断中。

概率与统计

- 随机变量和概率分布：用于建模和分析随机现象，如正态分布、泊松分布。
- 贝叶斯定理：用于更新概率分布，广泛应用于贝叶斯网络和贝叶斯推断。
- 最大似然估计（MLE）和最大后验估计（MAP）：用于参数估计和模型选择。
- 马尔可夫链和马尔可夫决策过程（MDP）：用于建模和解决决策过程，特别是在强化学习中。
- 蒙特卡罗方法：用于数值积分和概率分布的近似计算。

优化理论

- 凸优化：研究凸函数和凸集合的优化问题，如LASSO和支持向量机（SVM）。
- 非凸优化：处理复杂的优化问题，如深度学习中的神经网络训练。
- 梯度下降和随机梯度下降：用于迭代优化问题。
- 线性规划和整数规划：用于资源分配和组合优化问题。
- 拉格朗日乘数法：用于含约束条件的优化问题。

图论

- 图的表示：用于建模关系和网络，如邻接矩阵和邻接表。
- 最短路径算法：如Dijkstra算法，用于路径规划和网络分析。
- 图匹配和图划分：用于社交网络分析、图像分割和聚类。
- 网络流和最大流问题：用于优化运输和流量问题。

信息论

- 熵和互信息：用于衡量信息量和依赖关系，在特征选择和通信中应用。
- 编码理论：研究数据压缩和错误检测/纠正码，如香农编码和汉明码。
- 卡尔曼滤波：用于动态系统的状态估计和信号处理。

计算理论

- 图灵机和计算复杂性：研究计算问题的复杂性和可计算性，如P与NP问题。
- NP完全性：研究问题的求解难度和算法的效率，判定问题是否具有多项式时间复杂度解。
- 自动机理论：研究有限状态机和形式语言，应用于编译器和正则表达式。

离散数学

- 组合数学：用于分析和解决离散结构的问题，如组合优化、排列和组合。
- 布尔代数：用于逻辑推理和数理逻辑，特别是在逻辑编程和电路设计中。
- 图论和网络流：用于优化和分析网络结构。

动态系统与微分方程

- 常微分方程（ODEs）：用于建模连续时间的动态系统，如物理模拟和控制理论。
- 偏微分方程（PDEs）：用于建模空间和时间的连续变化，如图像处理和物理场模拟。
- 分岔理论和混沌理论：研究动态系统的稳定性和非线性行为。

几何与拓扑

- 微分几何：用于分析曲面和形状，如在计算机视觉和图像处理中的应用。
- 拓扑数据分析（TDA）：用于数据的形状分析和降维，应用于高维数据的可视化。
- 计算几何：研究几何对象的算法和数据结构，用于计算机图形学和机器人学。

变分法与计算几何

- 变分法：用于优化问题的解，如在图像处理、路径规划和物理模拟中的应用。
- 计算几何：研究几何对象的算法和数据结构，用于图形学、计算机视觉和CAD系统。

数值分析

- 数值微分和积分：用于逼近和求解微分方程和积分方程。
- 线性代数的数值方法：如QR分解、LU分解，用于求解线性方程组。
- 优化的数值方法：如牛顿法、共轭梯度法，用于高效求解优化问题。

笛卡尔几何与代数几何

- 笛卡尔几何：用于解析几何问题，如直线和曲线的表示与计算。
- 代数几何：研究多项式方程的解集结构，用于计算机代数系统和图形学。

逻辑与集合论

- 一阶逻辑和高阶逻辑：用于形式化推理和验证。
- 集合论：研究集合的性质和操作，是数学的基础理论之一。

\section{预备知识}

要学习深度学习，首先需要先掌握一些基本技能。所有机器学习方法都涉及从数据中提取信息。因此，我们先学习一些关于数据的实用技能，包括存储、操作和预处理数据。

机器学习通常需要处理大型数据集。我们可以将某些数据集视为一个表，其中表的行对应样本，列对应属性。线性代数为人们提供了一些用来处理表格数据的方法。我们不会太深究细节，而是将重点放在矩阵运算的基本原理及其实现上。

深度学习是关于优化的学习。对于一个带有参数的模型，我们想要找到其中能拟合数据的最好模型。在算法的每个步骤中，决定以何种方式调整参数需要一点微积分知识。本章将简要介绍这些知识。幸运的是，`autograd`包会自动计算微分，本章也将介绍它。

机器学习还涉及如何做出预测：给定观察到的信息，某些未知属性可能的值是多少？要在不确定的情况下进行严格的推断，我们需要借用概率语言。

最后，官方文档提供了本书之外的大量描述和示例。在本章的结尾，我们将展示如何在官方文档中查找所需信息。

本书对读者数学基础无过分要求，只要可以正确理解深度学习所需的数学知识即可。但这并不意味着本书中不涉及数学方面的内容，本章会快速介绍一些基本且常用的数学知识，以便读者能够理解书中的大部分数学内容。如果读者想要深入理解全部数学内容，可以进一步学习本书数学附录中给出的数学基础知识。

-2.1.数据操作
%-[2.1.1.入门](https：//zh-v2.d2l.ai/chapter_preliminaries/ndarray.html#id2)
%-[2.1.2.运算符](https：//zh-v2.d2l.ai/chapter_preliminaries/ndarray.html#id3)
%-[2.1.3.广播机制](https：//zh-v2.d2l.ai/chapter_preliminaries/ndarray.html#subsec-broadcasting)
%-[2.1.4.索引和切片](https：//zh-v2.d2l.ai/chapter_preliminaries/ndarray.html#id5)
%-[2.1.5.节省内存](https：//zh-v2.d2l.ai/chapter_preliminaries/ndarray.html#id6)
%-[2.1.6.转换为其他Python对象](https：//zh-v2.d2l.ai/chapter_preliminaries/ndarray.html#python)
%-[2.1.7.小结](https：//zh-v2.d2l.ai/chapter_preliminaries/ndarray.html#id7)
%-[2.1.8.练习](https：//zh-v2.d2l.ai/chapter_preliminaries/ndarray.html#id8)
%-2.2.数据预处理
%-[2.2.1.读取数据集](https：//zh-v2.d2l.ai/chapter_preliminaries/pandas.html#id2)
%-[2.2.2.处理缺失值](https：//zh-v2.d2l.ai/chapter_preliminaries/pandas.html#id3)
%-[2.2.3.转换为张量格式](https：//zh-v2.d2l.ai/chapter_preliminaries/pandas.html#id4)
%-[2.2.4.小结](https：//zh-v2.d2l.ai/chapter_preliminaries/pandas.html#id5)
%-[2.2.5.练习](https：//zh-v2.d2l.ai/chapter_preliminaries/pandas.html#id6)
%-2.3.线性代数
%-[2.3.1.标量](https：//zh-v2.d2l.ai/chapter_preliminaries/linear-algebra.html#id2)
%-[2.3.2.向量](https：//zh-v2.d2l.ai/chapter_preliminaries/linear-algebra.html#id3)
%-[2.3.3.矩阵](https：//zh-v2.d2l.ai/chapter_preliminaries/linear-algebra.html#id5)
%-[2.3.4.张量](https：//zh-v2.d2l.ai/chapter_preliminaries/linear-algebra.html#id6)
%-[2.3.5.张量算法的基本性质](https：//zh-v2.d2l.ai/chapter_preliminaries/linear-algebra.html#id7)
%-[2.3.6.降维](https：//zh-v2.d2l.ai/chapter_preliminaries/linear-algebra.html#subseq-lin-alg-reduction)
%-[2.3.7.点积（DotProduct）](https：//zh-v2.d2l.ai/chapter_preliminaries/linear-algebra.html#dot-product)
%-[2.3.8.矩阵-向量积](https：//zh-v2.d2l.ai/chapter_preliminaries/linear-algebra.html#id10)
%-[2.3.9.矩阵-矩阵乘法](https：//zh-v2.d2l.ai/chapter_preliminaries/linear-algebra.html#id11)
%-[2.3.10.范数](https：//zh-v2.d2l.ai/chapter_preliminaries/linear-algebra.html#subsec-lin-algebra-norms)
%-[2.3.11.关于线性代数的更多信息](https：//zh-v2.d2l.ai/chapter_preliminaries/linear-algebra.html#id14)
%-[2.3.12.小结](https：//zh-v2.d2l.ai/chapter_preliminaries/linear-algebra.html#id16)
%-[2.3.13.练习](https：//zh-v2.d2l.ai/chapter_preliminaries/linear-algebra.html#id17)
%-2.4.微积分
%-[2.4.1.导数和微分](https：//zh-v2.d2l.ai/chapter_preliminaries/calculus.html#id2)
%-[2.4.2.偏导数](https：//zh-v2.d2l.ai/chapter_preliminaries/calculus.html#id3)
%-[2.4.3.梯度](https：//zh-v2.d2l.ai/chapter_preliminaries/calculus.html#subsec-calculus-grad)
%-[2.4.4.链式法则](https：//zh-v2.d2l.ai/chapter_preliminaries/calculus.html#id5)
%-[2.4.5.小结](https：//zh-v2.d2l.ai/chapter_preliminaries/calculus.html#id6)
%-[2.4.6.练习](https：//zh-v2.d2l.ai/chapter_preliminaries/calculus.html#id7)
%-2.5.自动微分
%-[2.5.1.一个简单的例子](https：//zh-v2.d2l.ai/chapter_preliminaries/autograd.html#id2)
%-[2.5.2.非标量变量的反向传播](https：//zh-v2.d2l.ai/chapter_preliminaries/autograd.html#id3)
%-[2.5.3.分离计算](https：//zh-v2.d2l.ai/chapter_preliminaries/autograd.html#id4)
%-[2.5.4.Python控制流的梯度计算](https：//zh-v2.d2l.ai/chapter_preliminaries/autograd.html#python)
%-[2.5.5.小结](https：//zh-v2.d2l.ai/chapter_preliminaries/autograd.html#id5)
%-[2.5.6.练习](https：//zh-v2.d2l.ai/chapter_preliminaries/autograd.html#id6)
%-2.6.概率
%-[2.6.1.基本概率论](https：//zh-v2.d2l.ai/chapter_preliminaries/probability.html#id2)
%-[2.6.2.处理多个随机变量](https：//zh-v2.d2l.ai/chapter_preliminaries/probability.html#id5)
%-[2.6.3.期望和方差](https：//zh-v2.d2l.ai/chapter_preliminaries/probability.html#id12)
%-[2.6.4.小结](https：//zh-v2.d2l.ai/chapter_preliminaries/probability.html#id13)
%-[2.6.5.练习](https：//zh-v2.d2l.ai/chapter_preliminaries/probability.html#id14)
%-2.7.查阅文档
%-[2.7.1.查找模块中的所有函数和类](https：//zh-v2.d2l.ai/chapter_preliminaries/lookup-api.html#id2)
%-[2.7.2.查找特定函数和类的用法](https：//zh-v2.d2l.ai/chapter_preliminaries/lookup-api.html#id3)
%-[2.7.3.小结](https：//zh-v2.d2l.ai/chapter_preliminaries/lookup-api.html#id8)
%-[2.7.4.练习](https：//zh-v2.d2l.ai/chapter_preliminaries/lookup-api.html#id9)


\section{浅谈图像处理：从矩阵到深度学习}

目  录
摘要
ABSTRACT
第1章  引言	1
1.1 图像识别如何作用于我们的生活？	1
1.2 追根溯源——其中的数学方法	1
1.3 报告主体结构与主要内容	1
第2章  从矩阵到深度学习	2
2.1 矩阵	2
2.1.1 矩阵的起源	2
2.1.2 矩阵的发展	2
2.1.3 矩阵的基本操作	2
2.2 传统图像处理	5
2.2.1基础知识——图像矩阵化	5
2.2.2 传统图像处理的起源	10
2.2.3 传统图像处理的发展	12
2.2.4 传统图像处理的应用	19
2.3 深度学习与图像处理——三种模型	20
2.3.1 起源	20
2.3.2 单层感知器模型	21
2.3.3 多层感知器模型	21
2.3.4 卷积神经网络模型	22
2.4 深度学习在图像处理上的发展——以目标检测为例	26
2.4.1 起源	26
2.4.2 R-CNN	27
2.4.3 SPP-Net	30
2.4.4 Fast R-CNN	31
2.4.5 Faster R-CNN	32
2.4.6 FPN	34
2.4.7 单阶段算法与两阶段算法	35
2.4.8 总结	35
2.4.9 深度学习图像处理举例	35
第3章  体会启发与结语	38
3.1体会与启发	38
3.2 结语	39
参考文献	40
致  谢	41

摘要
本论文从代数学矩阵开始，讲述了将图像矩阵（数字）化的重要意义，以及如何使用矩阵分析与变换的方法来分析与修改图像，并介绍了传统图像处理技术的建立与发展过程，若干重要算法以及一些弊端，最后将图像处理技术指向了深度学习与神经网络模型，引入人工智能协助以达到更先进的技术水平，并展示了一个具体实例。

\section{第1章  引言}
\subsection{1.1 图像识别如何作用于我们的生活？}

	图像识别是计算机视觉和人工智能领域的重要组成部分，其终极目标是使计算机具有分析和理解图像内容的能力。图像识别是一个综合性的问题，涵盖图像匹配、图像分类、图像检索、人脸检测、行人检测等技术，并在互联网搜索引擎、自动驾驶、医学分析、遥感分析等领域具有广泛的应用价值。
具体例子：最近，由于疫情封锁，室内运动成了当下人们主要的运动方式，一款用于督促人们运动的app应运而生，那就是天天跳绳，而它就是通过识别人的肢体动作图像来对人的运动做出判断。这样的例子数不胜数。

\subsection{1.2 追根溯源——其中的数学方法}
	将图像进行数字（矩阵）化后，我们就可以使用对于矩阵的分析方法来分析图像了，由矩阵的一些基本运算即可达到修改图像的目的，例如可以使用仿射变换来对图像做平移、放大、旋转，剪切等操作。更可以引入更加深入的算法和数学知识来进行一些更为先进的图像处理方式，例如使用高等数学中的卷积概念进行图像分析，使用Canny算法实现边缘检测，以及借许多数学知识与编程思维使用深度学习技术训练，达到使用人工智能协助图像分析的目的。

\subsection{1.3 报告主体结构与主要内容}
	本报告从矩阵知识引入，介绍了其起源（2.1.1）、发展（2.1.2）、基础知识（2.1.3），这一部分由王晨阳同学负责。
然后，我们探究了其在传统图像处理中的应用——图像矩阵化（2.2），层层深入，揭示了现代图像处理中矩阵的作用，这一部分由谭诚同学负责。
在传统图像处理（2.2）中我们探究了其方法（2.2.1），起源（2.2.2），发展（2.2.3）及应用（2.2.4），这一部分由汪坤灿同学负责。
在深度学习与图像处理中，我们介绍了若干种深度学习的神经网络模型，以及若干个深度学习的算法，这一部分由史晓磊同学负责。最后，我们使用Python-OpenCV库以及调用了Github上一个已经训练好的模型，完成了一次对图片中人脸的识别检测，这一部分由谭诚同学负责。
小组中，我们各自安排撰写各自的文档、ppt，以及录制视频中的各个部分。最后由谭诚同学整合、排版，以及剪辑。
	 
\section{第2章  从矩阵到深度学习}

去年的某一天，一张利用AI修复的林徽因的照片冲上了热搜。可以看出，右边由AI修复的照片清晰度有了巨大的提升，尽管仍有些模糊，但是对于照片的关键部分--人脸已经有了巨大的提升。可以看出后者非常像网红脸。

那么人工智能是如何修复照片的呢？那就是对大量人脸数据的特征进行提取。但是对这么大量的数据如何分析呢？

如上图，在电路中，该图是最简单的一个电路模型。其输出端Y的值即为输入端Y的值，其值有两种，0和1。因此，我们可以通俗的认为，0代表否，1代表是。因此，对于一个特征的有无，0代表无，1代表有。这样就可以用数字记录一个事物的特征。同样的，利用一串数组，可以记录一系列的特征。那么利用矩阵，就可以记录大量的特征。
本报告将涉及数学概念中的“矩阵（Martix）”。

\subsection{2.1 矩阵}

在数学中，矩阵（Matrix）是一个按照长方阵列排列的复数或实数集合 [1]  ，最早来自于方程组的系数及常数所构成的方阵。这一概念由19世纪英国数学家凯利首先提出。作为解决线性方程的工具，矩阵也有不短的历史。成书最早在东汉前期的《九章算术》中，用分离系数法表示线性方程组，得到了其增广矩阵。在消元过程中，使用的把某行乘以某一非零实数、从某行中减去另一行等运算技巧，相当于矩阵的初等变换。但那时并没有现今理解的矩阵概念，虽然它与现有的矩阵形式上相同，但在当时只是作为线性方程组的标准表示与处理方式。

{2.1.1 矩阵的起源}

作为解决线性方程的工具，矩阵也有不短的历史。成书最迟在东汉前期的《九章算术》中，用分离系数法表示线性方程组，得到了其增广矩阵。在消元过程中，使用的把某行乘以某一非零实数、从某行中减去另一行等运算技巧，相当于矩阵的初等变换。但那时并没有现今理解的矩阵概念，虽然它与现有的矩阵形式上相同，但在当时只是作为线性方程组的标准表示与处理方式。
矩阵正式作为数学中的研究对象出现，则是在行列式的研究发 展起来后。逻辑上，矩阵的概念先于行列式，但在实际的历史上则恰好相反。日本数学家关孝和（1683年）与微积分的发现者之一戈特弗里德·威廉·莱布尼茨（1693年）近乎同时地独立建立了行列式论。其后行列式作为解线性方程组的工具逐步发展。1750年，加布里尔·克拉默发现了克莱姆法则 。

	矩阵最早出现在我国的九章算术中，在九章算术方程篇中，刘徽就给出过“方程”的定义，不过他所定义的方程，指的是由线性方程得到的数码矩阵，其本质就是线性方程组的增广矩阵，不仅如此，他还提出过解线性方程组的两种方法，即利用该方阵进行“直除”和“互乘”运算，“直除”就是将某行的同一倍数加到另一行上，“互乘”则是以一个不为零的数乘该方阵的某一行上，这里已经涉及到了矩阵初等变换的知识。
	
{2.1.2 矩阵的发展}
	1801年，德国数学家高斯首次将线性方程的全部系数放在一起，并提出了高斯消元法（这种方法在今天的计算机解线性方程组的问题中仍在使用），给出了矩阵的雏形；1844年，德国数学家爱森斯坦首次讨论了变换（也就是今天的“矩阵”）的乘积，为此后的数学家打开了一条新路。
1850年，英国数学家西尔维斯特首先使用矩阵一词；1858年，英国数学家凯莱发表《关于矩阵理论的研究报告》。他首先将矩阵作为一个独立的数学对象加以研究，并在这个主题上首先发表了一系列文章，因而被认为是矩阵论的创立者，他给出了现在通用的一系列定义，如两矩阵相等、零矩阵、单位矩阵、两矩阵的和、一个数与一个矩阵的数量积、两个矩阵的积、矩阵的逆、转置矩阵等，并且凯莱还注意到矩阵的乘法是可结合的，但一般不可交换，且m*n矩阵只能用n*k 矩阵去右乘。1854 年，法国数学家埃米尔特使用了"正交矩阵"这一术语，但他的正式定义直到1878年才由德国数学家费罗贝尼乌斯发表。1879年，费罗贝尼乌斯引入秩的概念。至此，矩阵的体系基本上完全建立起来了。


      矩阵的现代概念在19世纪逐渐形成。1800年代，高斯和威廉·若尔当建立了高斯—若尔当消去法。1844年，德国数学家费迪南·艾森斯坦（F.Eisenstein）讨论了“变换”（矩阵）及其乘积。1850年，英国数学家詹姆斯·约瑟夫·西尔维斯特（James Joseph Sylvester）首先使用矩阵一词 。
英国数学家凯利被公认为矩阵论的奠基人。他开始将矩阵作为独立的数学对象研究时，许多与矩阵有关的性质已经在行列式的研究中被发现了，这也使得凯利认为矩阵的引进是十分自然的。他说：“我决然不是通过四元数而获得矩阵概念的；它或是直接从行列式的概念而来，或是作为一个表达线性方程组的方便方法而来的。”他从1858年开始，发表了《矩阵论的研究报告》等一系列关于矩阵的专门论文，研究了矩阵的运算律、矩阵的逆以及转置和特征多项式方程。凯利还提出了凯莱-哈密尔顿定理，并验证了3×3矩阵的情况，又说进一步的证明是不必要的。哈密尔顿证明了4×4矩阵的情况，而一般情况下的证明是德国数学家弗罗贝尼乌斯（F.G.Frohenius）于1898年给出的 。
1854年时法国数学家埃尔米特（C.Hermite）使用了“正交矩阵”这一术语，但他的正式定义直到1878年才由费罗贝尼乌斯发表。1879年，费罗贝尼乌斯引入矩阵秩的概念。至此，矩阵的体系基本上建立起来了。
    无限维矩阵的研究始于1884年。庞加莱在两篇不严谨地使用了无限维矩阵和行列式理论的文章后开始了对这一方面的专门研究。1906年，希尔伯特引入无限二次型（相当于无限维矩阵）对积分方程进行研究，极大地促进了无限维矩阵的研究。在此基础上，施密茨、赫林格和特普利茨发展出算子理论，而无限维矩阵成为了研究函数空间算子的有力工具 。


\subsection{矩阵的应用现状}
矩阵在机器学习和深度学习上面的应用，有以下几个。
1.	在将数据集应用于机器学习模型之前，需要对数据集做预处理，比如给数据去噪、降维，这样子可以减少计算资源和提高模型预测准确度。常用的去噪降维的方法是主成分分析（PCA），里面就用到了矩阵，具体步骤为（1）求解数据集的协方差矩阵的特征值和特征向量，（2）根据特征值从大到小排序（3）选取前面几个特征值对应的特征向量作为新的坐标基向量。（4）将所有的数据集映射到这个新的坐标系中，从而得到降噪降维之后的数据集了。在第（3）步中只选取前面最大的几个特征值所对应的特征向量，是因为小特征值所蕴含的信息量比较小，它们很可能是一些噪声数据。
2.	另一个在机器学习中的应用是奇异值分解，奇异值分解的公式为M=U∑V  , 在Netflix Prize电影推荐系统中，根据用户对电影的评分来建模的推荐系统使用到了奇异值分解的思想。
3. 在深度学习模型的框架中，数据是以矩阵的形式组织的，比如一次输入64张彩色图片，输入数据是64x3x32x32 维度的。这样子一方面可以简化表示，另一方面可以提高计算速度，特别是GPU能够针对这个进行优化，如果不以矩阵的形式表示数据，那么GPU的优势就没发挥出来。

矩阵是人工智能基础范畴中不可或缺的一部分。
必备的数学知识是理解人工智能不可或缺的要素，所有的人工智能技术归根到底都建立在数学模型之上，而这些数学模型又都离不开线性代数的理论框架。
线性代数不仅仅是人工智能的基础，更是现代数学和以现代数学作为主要分析方法的众多学科的基础。从量子力学到图像处理都离不开向量和矩阵的使用。而在向量和矩阵背后，线性代数的核心意义在于提供了⼀种看待世界的抽象视角：万事万物都可以被抽象成某些特征的组合，并在由预置规则定义的框架之下以静态和动态的方式加以观察。线性代数是 AI 专家必须掌握的知识.线性代数是一套非常通用的思想和工具，可以应用于各个领域。
如果你只想把人工智能和机器学习的工具当作一个黑匣子，那么你只需要足够的数学计算就可以确定你的问题是否符合模型使用。如果你想提出新想法，线性代数则是你必须要学习的东西。并不是说你需要学习有关数学的所有知识，这样会耽搁于此，失去研究其他更重要的东西（如微积分/统计）的动力。
线性代数是应用数学领域，这是人工智能专家离不开的。如果你不精通这个领域，你永远不会成为一个好的AI专家。正如斯凯勒·斯皮克曼所说：“线性代数是21世纪的数学。”线性代数有助于产生新的想法，这就是为什么它是人工智能科学家和研究人员必须学习的东西。他们可以用标量、向量、张量、矩阵、集合和序列、拓扑、博弈论、图论、函数、线性变换、特征值和特征向量的概念来建立模型。
向量 矩阵理论
在科幻电影中，你通常会看到，通过执行一些类似于神经系统的计算结构，产生了一个神经网络，它生成神经元之间的连接，以匹配人类大脑的推理方式。将矩阵的概念引入到神经网络的研究中
神经网络可以通过形成三层人工神经元来实现非线性假设：
1. 输入层   2. 隐藏层   3.输出层
特征向量
搜索引擎排名的科学是建立在数学科学的基础上的。页面排名，这正是谷歌作为一家基于数学视角的公司的基础。页面排名是一种算法，最初由拉里·佩奇和谢尔盖·布林在他们的研究论文《大型超文本Web搜索引擎的剖析》中提出。在这一巨大突破的背后，使用了数百年来众所周知的特征值和特征向量的基本概念。
机器人爬虫程序首先检索网页，然后通过分配页面排名值对其进行索引和编目。每个页面的可信度取决于该页面的链接数。


\subsection{2.1.3 矩阵的基本操作}

2.1.3.1 加减法

对应位置相加减。例如：
%[(1&0@0&1)]+[■(0&1@1&0)]=[■(1&1@1&1)]
%[■(1&0@0&1)]-[■(0&1@1&0)]=[■(1&-1@-1&1)]
需要注意的是，这两个矩阵必须行列数均相同。

2.1.3.2 矩阵数乘

矩阵的所有元素都乘上这个数即可。
%$k⋅[■(0&1@1&0)]=[■(0&k@k&0)]$

2.1.3.3 矩阵乘法

两个矩阵的乘法仅当第一个矩阵A的列数和另一个矩阵B的行数相等时才能定义。如果矩阵A是m×n矩阵（m行n列），矩阵B是n×k矩阵（n行k列），则它们的乘积C是一个m×k矩阵。对于矩阵C的每个元素，它如下计算：
% c_(i,j)=a_(i,1) b_(1,j)+a_(i,2) b_(2,j)+⋯+a_(i,n) b_(n,j)=∑_(k=1)^n▒〖a_(i,k) b_(k,j) 〗
%[■(1&0@0&1)]⋅[■(0&1@1&0)]=[■(0&1@1&0)]
容易看出，单位矩阵乘任何矩阵，不论左乘还是右乘，都不会改变矩阵。这在传统的线性代数教材中已经被证明很多次了，这里不再赘述。

2.1.3.4 矩阵转置

将每个元素的行数和列数互换即可。
%[■(1&2@3&4)]^T=[■(1&3@2&4)]

2.1.3.5 共轭矩阵

共轭是复数域的概念，只是把矩阵中每个元素取其共轭的复数即可。
%A=[■(3+i&5@2-2i&i)],    A ̅=[■(3-i&5@2+2i&-i)]

2.1.3.6 矩阵的特征值及特征向量

设A是n阶方阵，如果数λ和n维非零列向量x使关系式Ax=λx成立，那么这样的数λ称为矩阵A特征值，非零向量x称为A的对应于特征值λ的特征向量。对于矩阵的特征值及特征向量，在传统线性代数教材中已经有了足够充分的解释，并对其性质和计算方法也有了足够多的介绍，这里不再赘述了。

矩阵的特征值和特征向量可以揭示线性变换的深层特性 ，我们稍后会展开介绍其具体应用，并简述其实现原理。

2.1.3.7 线性变换

在数学上，如果满足下式，那么映射F(a)就是线性的：

F(a + b) = F(a) + F(b) 以及  F(ka) = kF(a)

如果映射F保持了基本运算：加法和数量乘，那么就可以称该映射为线性的。在这种情况下，将两个向量相加然后再进行变换得到的结果和先分别进行变换再将变换后的向量相加得到的结果相同。同样，将一个向量数量乘再进行变换和先进行变换再数量乘的结果也是一样的。

这个线性变换的定义有两条重要的引理：

(1) 映射F(a) = aM，当M为任意方阵时，此映射是一个线性变换，这是因为：
F(a + b)= (a + b)M = aM + bM = F(a)+ F(b)
F(ka) = (ka)M = k(aM) = kF(a)

(2) 零向量的任意线性变换的结果仍然是零向量。（如果F(0) = a，a ≠ 0。那么F不可能是线性变换。因为F(k0) = a，但F(k0) ≠ kF(0)，因此线性变换不会导致平移（原点位置上不会变化）。

2.1.3.8 矩阵的仿射变换

仿射变换是指线性变换后接着平移。因此仿射变换的集合是线性变换的超集，任何线性变换都是仿射变换，但不是所有仿射变换都是线性变换。

任何具有形式$ v^\prime  = vM + b$ 的变换都是仿射变换。

2.1.3.9 矩阵的可逆变换
如果存在一个逆变换可以“撤销”原变换，那么该变换是可逆的。换句话说，如果存在逆变换G，使得G(F(a)) = a，对于任意a，映射F(a)是可逆的。

存在非仿射变换的可逆变换，但暂不考虑它们。现在，我们集中精力于检测一个仿射变换是否可逆。一个仿射变换就是一个线性变换加上平移，显然，可以用相反的量”撤销“平移部分，所以问题变为一个线性变换是否可逆。

显然，除了投影以外，其他变换都能“撤销”。当物体被投影时，某一维有用的信息被抛弃了，而这些信息时不可能恢复的。因此，
所有基本变换除了投影都是可逆的。因为任意线性变换都能表达为矩阵，所以求逆变换等价于求矩阵的逆。如果矩阵是奇异的，则变换不可逆；可逆矩阵的行列式不为0。

2.1.3.10 矩阵的等角变换

如果变换前后两向量夹角的大小和方向都不改变，该变换是等角的。只有平移，旋转和均匀缩放是等角变换。等角变换将会保持比例不变，镜像并不是等角变换，因为尽管两向量夹角的大小不变，但夹角的方向改变了。所有等角变换都是仿射和可逆的。

2.1.3.11 矩阵的正交变换

术语“正交”用来描述具有某种性质的矩阵。正交变换的基本思想是轴保持互相垂直，而且不进行缩放变换。
平移、旋转和镜像是仅有的正交变换。长度、角度、面积和体积都保持不变。（尽管如此，但因为镜像变换被认为是正交变换，所
以一定要密切注意角度、面积和体积的准确定义）。
正交矩阵的行列式为1或者负1，所有正交矩阵都是仿射和可逆的。

2.2 传统图像处理

2.2.1 基础知识——图像矩阵化

2.2.1.1 灰度[9][10][11]
	在电子计算机领域中，灰度（Gray scale）数字图像是每个像素只有一个采样颜色的图像。这类图像通常显示为从最暗黑色到最亮的白色的灰度，尽管理论上这个采样可以是任何颜色的不同深浅，甚至可以是不同亮度上的不同颜色。灰度图像与黑白图像不同，在计算机图像领域中黑白图像只有黑白两种颜色，灰度图像在黑色与白色之间还有许多级的颜色深度。但是，在数字图像领域之外，“黑白图像”也表示“灰度图像”，例如灰度的照片通常叫做“黑白照片”。在一些关于数字图像的文章中单色图像等同于灰度图像，在另外一些文章中又等同于黑白图像。[12]

图2.2.1.1.1 原始彩色图片与灰度图

灰度图像经常是在单个电磁波频谱如可见光内测量每个像素的亮度得到的。用于显示的灰度图像通常用每个采样像素8 bits的非线性尺度来保存，这样可以有256种灰度（8bits就是2的8次方=256）。这种精度刚刚能够避免可见的条带有损，并且非常易于编程。在医学图像与遥感图像这些技术应用中经常采用更多的级数以充分利用每个采样10或12 bits的传感器精度，并且避免计算时的近似误差。在这样的应用领域流行使用16 bits即65536个组合（或65536种颜色）。[12][13]

2.2.1.2 RGB颜色系统[7]
	三原色光模式（RGB color model），又称RGB颜色模型或红绿蓝颜色模型，是一种加色模型，将红（Red）、绿（Green）、蓝（Blue）三原色的色光以不同的比例相加，以合成产生各种色彩光。
RGB颜色模型的主要目的是在电子系统中检测，表示和显示图像，比如电视和电脑，利用大脑强制视觉生理模糊化(失焦),将红绿蓝三原色子象素合成为一色彩象素,产生感知色彩(其实此真彩色并非加色法所产生的合成色彩，原因为该三原色光从来没有重叠在一起，衹是人类为了“想”看到色彩，大脑强制眼睛失焦而形成。红绿蓝三色模型在传统摄影中也有应用。在电子时代之前，基于人类对颜色的感知，RGB颜色模型已经有了坚实的理论支撑。
RGB是一种依赖于设备的颜色空间：不同设备对特定RGB值的检测和重现都不一样，因为颜色物质（荧光剂或者染料）和它们对红、绿和蓝的单独响应水平随着制造商的不同而不同，甚至是同样的设备不同的时间也不同。[7]
     
图2.2.1.1.1 RGB颜色空间示意图

一个颜色显示的描述是由三个数值控制的，他分别为R、G、B。当三个数值为最大（255）时，显示为白色，当三个数值最小（0）时，显示为黑色。数值表示可以使用以下几种不同的方式：从0到1之间可用的浮点数来表示、从0 \%到100 \%用百分比表示，或者使用0到255之间的整数以八位数字表示，通常表示为十进制和十六进制的数值。[7]

2.2.1.3 Gamma校正[8][11]
	
	伽马校正（Gamma correction） 又叫伽马非线性化（gamma nonlinearity）、伽马编码（gamma encoding） 或是就只单纯叫伽马（gamma）。是用来针对影片或是影像系统里对于光线的辉度（luminance）或是三色刺激值（tristimulus values）所进行非线性的运算或反运算。最简单的例子中，伽马校正是由下幂定律公式所定义的：
$V_out=AV_in^\gamma$
	
	其中A是一个常量，输入和输出的值都为非负实数值。一般地来说在A=1的通常情况下，输入输出的值的范围都是在0到1之间。伽马值~$\gamma < 1$的情况有时被称作编码伽马值（encoding gamma），而执行这个编码运算所使用上述幂定律的过程也叫做伽马压缩（gamma compression）；相对地，伽马值~$\gamma > 1$的情况有时也被称作解码伽马值（decoding gamma），而执行这个解码运算所使用上述幂定律的过程也叫做“伽马展开（gamma expansion）”。[8][11]
	
	RGB值与功率并非简单的线性关系，而是幂函数关系，这个函数的指数称为Gamma值，一般为2.2，而这个换算过程，称为Gamma校正。为什么显示器要Gamma校正呢？因为人眼对亮度的感知和物理功率不是成正比的，而是幂函数的关系，这个函数的指数通常为2.2，称为Gamma值。所以RGB中的灰度值，为了考虑到较小的存储范围(0~255)和较平衡的亮暗部比例，所以需要进行Gamma校正，而不是直接对应功率值，因此RGB值RGB颜色值不能简单直接相加，而是必须用2.2次方换算成物理光功率后才能进行下一步计算。
	
	由于Gamma校正，在计算机显示设备上的颜色输出的强度通常不是直接正比于在图像文件中R, G和B值。就是说，即使值0.5非常接近于0到1.0（完全强度）的一半，计算机显示器在显示 (0.5, 0.5, 0.5)时候的光强度通常（在标准2.2-gamma CRT/LCD上）是在显示 (1.0, 1.0, 1.0)时候的大约22 \%，而不是50 \%。

图2.2.1.3.1 Gamma校正
	
	为图像进行伽马编码的目的是用来对人类视觉的特性进行补偿，从而根据人类对光线或者黑白的感知，最大化地利用表示黑白的数据位或带宽。在通常的照明（既不是漆黑一片，也不是令人目眩的明亮）的情况下，人类的视觉大体有伽马或者是幂函数的性质。如果不将图像进行伽马编码，那么数据位或者带宽的利用就会分布不均匀——会有过多的数据位或者带宽用来表示人类根本无法察觉到的差异，而用于表示人类非常敏感的视觉感知范围的数据位或者带宽又会不足。图像的伽马编码并不是必须的（甚至有的时候会适得其反），浮点数格式的颜色值已经提供了一部分对数曲线的线性估计。[8][11]

2.2.1.4 图像的数字（矩阵）化

	位图[14][15]（英语：Bitmap，台湾称为点阵图），又称栅格图（Raster graphics），是使用像素阵列(Pixel-array/Dot-matrix点阵)来表示的图像。我们这篇课题所涉及的图像，都是在指位图，而不是矢量图。
	我们先来考虑灰度图在OpenCV中如何存储。OpenCV是一个跨平台计算机视觉和机器学习软件库，我们稍后会介绍它。这里我们会用简单的代码来读取一个图片并观察它的输出结果，来推断出图像是如何数字化的。

C++代码
	以及我们需要读取的“1.jpg”图片：

“1.jpg” 图片（一个方格示意为一个像素）
	输出结果是：




事实上，这也就是矩阵：
%[■(0&255&0@255&0&255@0&255&0)]

	结合我们上面介绍过的RGB与灰度的知识，不难发现出：由于灰度仅由一个大小为0~255的八位无符号整数表示，所以每个像素点只需要用一个数来表示就可以了（黑色的灰度是0，白色的灰度是255）.

	接下来，我们再来考虑RGB色彩系统下的彩色图像的数字（矩阵）化。








C++代码
	以及我们需要读取的“1.jpg”图片：

“1.jpg” 图片（一个方格示意为一个像素）
	输出结果是：




	事实上，这也就是矩阵：
%[■(0&0&255@0&255&0@255&0&0)■(      255&0&0@      0&0&255@      0&255&0)■(      0&0&255@      0&255&0@      255&0&0)]

	我们在”1.jpg”中预设的三种颜色都是三原色，也就是：以(R,G,B)三元组表示一个颜色的话，红色为(255,0,0)，绿色为(0,255,0)，蓝色为(0,0,255)。

	而观察我们的矩阵，这是一个3x3像素的图像，而这个矩阵却是一个3x9矩阵，结合我们上面介绍的RGB颜色系统，可以得知一个颜色被三个矩阵元素表示，且这三个元素分别就是RGB三个值中的某个值。

	然而，我们图像中第一行的颜色顺序为“红-蓝-红”，按照(R,G,B)三元组表示，它们应该为”(255,0,0) - (0,0,255) - (255,0,0)”，然而矩阵中却是”(0,0,255) - (255,0,0) - (0,0,255)”，这说明实际在OpenCV中存储时的顺序为(B,G,R)。

	到这里为止，所有的位图都可以如上数字化了。在RGB颜色系统中，我们还可以分析一个图像中各原色通道的占比，并用灰度图输出表示：

	例如，上图橙子果肉的颜色用(R,G,B)三元组表示大概为(168,102,0)，在red通道图中它最明亮，在green通道图中其次，在blue通道图中最暗。在灰度中，255为白色，0为黑色，这表示出了图像中各个位置的原色比例。

2.2.2 传统图像处理的起源[1]

	数字图像处理，是现代信息处理形式的代表，具有画面处理程序一体化、图像处理清晰度高等特性。较之传统的图像处理，如胶片剪辑、胶片冲洗等，数字图像处理利用计算机对图像进行加工与处理，满足了图片数字化、可数码操作的需要。1921年[1]，人们通过海底电缆将图像由伦敦传输至纽约，标志着数字图像处理系统的诞生。首个数字图像处理系统名为“巴特兰”（Bartlane），该系统使用5个不同的灰度级来编码图像。而到了1929年，“巴特兰”编码图像的这一 能力已经扩展到了15级。

图2.2.1 1921年经“巴特兰”系统处理的数字图片

图2.2.2 1929年经“巴特兰”系统处理的数字图片

20世纪50年代，人们开始 对数字图像处理技术进行系统的研究。1964年，美国加州的喷气推进实验室使用数字计算机处理了“徘徊者7号”月球探测器发回的宇宙照片。

图2.2.3 “徘徊者7号”月球探测器
该事件为数字图像处理的一个重要里程碑，标志着第三代计算机问世后数字图像处理概念开始得到运用。此后数字图像处理技术发展迅速并于不久之后形成一门独立的学科，初步建立了起来。

图2.2.4 “徘徊者7号”月球探测器传回的照片

2.2.3 传统图像处理的发展

	在数字图像处理成为一门独立的学科之后的三四十年，随着各领域对图像处理的要求越来越高，数字图像处理技术得到了更加深入、广泛和迅速的发展。国内外学者在图像处理的各个方面取得了令人瞩目的成果……

	伴随着数字图像处理的不断发展[2]，各种处理技术层出不穷如图像数字化、图像变换、图像增强、图像恢复、图像数据压缩、图像边缘检测、图像分割、图像特征分析、图像配准、图像融合、图像分类、图像识别、基于内容的图像检索、图像数字水印等……

	下面将为大家以传统数字图像处理技术中的图像的增强、图像的滤波、图像的边缘检测的典型算法为例，为大家展示下数字图像处理技术的魅力。
	
2.2.3.1 图像增强算法[2]

图像增强算法用于增强对所需信息的辨别能力但同时也可能损失一些其他信息。因此在实际的应用中要对图像进行多种变换，然后再从中选取效果较好的一个。下面我们将以图像增强算法中的灰度变换为例来解释说明图像增强的一种实现。

图2.2.5 利用图像增强算法突出水中鱼类

灰度变换：
	灰度变换是指根据某种目标条件来确立变换关系逐然后逐点改变图像中每一个像素的灰度值。[3]让我们假设原图像像素的灰度值D=f(x,y)，处理后图像像素的灰度值记为$Ｄ\prime＝ｇ(x,y)$,则灰度增强可表示为：$D^\prime=T(D)$其中 D和 $D^\prime$都在图像的灰度范围之内。函数T(D)称为灰度变换函数，它描述了输入灰度值和输出灰度值之间的转换关系。 当我们确定了灰度变换函数后，一个具体的灰度增强方法便确立了。根据T(D)表达式的不同将灰度变换分为线性变换和非线性变换。
	其中的线性变换如下所示：
	若 ｇ(x,y)＝Ｔ[f(x,y)]是一个线性或分段线性的单值函数例如：
	对某个定义域内的灰度 g(x,y)=T[f(x,y)]＝a*f(x,y)+b，则由它确定的灰度变换称为灰度线性变换简称线性变换。式中参数 ａ为线性函数的斜率ｂ为 线性函数的在 y轴的截距。
	
线性灰度变换具有算法较易实现，使用方便等特性。但是其可扩展性较差。

一般地，我们更多会使用非线性灰度变换。


图2.2.6 对钙化点的医学影像使用线性灰度变换处理前后对比图

相对线性变换，当映射关系为非线性函数时，该变换显然为非线性变换。常用的非线性变换有对数变换和指数变换。
其中的对数变换如右所示：g(x,y) =c*ln[f(x,y)+1] 
公式中的 c是为了调整曲线的位置和形状而引入的参数。由于对数函数的自身特性，当我们想要对图像的低灰度区进行较大的扩展而对高灰度进区行压缩时可采取此变换。此处的灰度对数变换能使图像灰度分布与人的视觉特性相匹配。
其中的指数变换如右所示：$g(x,y) ＝b^ {c*[f(x,y)-a] －1}$
公式中a、b、c是为了调整曲线的位置和形状。由于指数函数的特性，故该变换能对图像的高灰度区域给予较大的扩展。

	综上，我们可以得出，当我们关注的灰度区不同时，我们可以选用不同的灰度变换来达到理想的效果。

图2.2.7 对钙化点的医学影像使用非线性灰度变换处理前后对比图

2.2.3.2 图像滤波算法

众所周知，由于外部或内部的各种原因， 数字图像不可避免地会有一些噪声[4]，如拍摄图片时光照不足会产生噪点，扫描图片时我们也无法得到百分之百不含杂质的图像。这便对后续的图像处理造成了影响。因此，图像去噪是数字图像处理工作中一个必不可少的部分，对于后续的各项其他处理都有重要意义。 应用图像滤波算法去噪已被应用于各项研究，涉及到了医疗勘探、雷达、文字等方面。发展至今，优化滤波算法，改进滤波算法的设计仍具有重要的意义。下面我们将以滤波算法中的中值滤波为例来说明图像滤波算法的一种实现。

图2.2.8 利用一种图像滤波算法来去除桥梁探伤图片的噪声
常规中值滤波：
不同于均值滤波中对像素矩阵中（x,y）邻域进行平均化的处理，中值滤波是一种非线性的处理方法。[5]
	  常规中值滤波首先用3×3大小的滑动窗口在彩色图像上自左向右、自上向下移动，设Sxy代 表以点(x,y)为中心、大小是3×3的矩形图像邻域像素值（窗口）的一组坐标：

再对Sxy中9个点的像素值（R,G,B）分别进行排序，一般使用冒泡排序法，得到三个序列Rij,Gij,Bij，以Rij 为例：

其中，f1≤f2≤f3≤…≤f9。之后，使用序列Rij的中间值替代原窗口中心点 (x,y)的R值f(x,y)。G值、B值操作过程相同，分别对（R,G,B）值完成滤波处理后，点(x,y)处的中值滤波处理完成，窗口移动到下一个像素点。

图2.2.9 常规中值滤波算法流程图

图2.2.10 利用常规中值滤波算法对加入脉冲噪声后的Lena图像处理前后对比图
通过对比图可以看到，常规中值滤波算法处理后虽然去除了大部分噪声但是同时也损失了不少图像细节，如Lena图中其脸颊和肩膀部分都损失了较多细节，通俗来说，图像某些部分变得较为模糊，仍有很大的改进空间。
自适应中值滤波：
自适应中值滤波算法是常规中值滤波算法的改进算法，在进行中值滤波处理之前，它会先判断采样点是否是噪声点，如果是噪声点，则进行滤波处理，如果不是噪声，则跳过当前采样点，移动到下一个采样点。
	和常规中值滤波算法的前半部分相同。（右图为自适应中值滤波算法流程图）首先用3×3大小的滑动窗口在彩色图像上自左向右、自上向下移动，设Sxy代 表以点(x,y)为中心、大小是3×3的矩形图像邻域像素值（窗口）的一组坐标：

再对Sxy中9个点的像素值（R,G,B）分别进行排序，一 般使用冒泡排序法，得到三个序列Rij,Gij,Bij，以Rij 为例：

其中，f1≤f2≤f3≤…≤f9。当排序后，中心点(x,y)的R值f(x,y)位于向量Rij（已排序） 的两端时，中心点(x,y)可能是一个噪声点。如果它同 时满足f(x,y)≤Tmin或f(x,y)≥Tmax，则中心点(x,y)判定为噪声点。参数Tmin和Tmax根据情况设定，一般Tmin取 40， Tmax取190比较合适。判定为噪声后，使用序列Rij的中间值f5替代原窗口中心点(x,y)的值f(x,y)，完成点(x,y)处的中值滤波处理，然后，窗口移动到下一个像素点，如此循环操作。可以看到自适应中值滤波算法在去除噪声的同时，图像细节损失较少。在大多数情况下，自适应中值滤波算法的处理效果是要优于常规中值滤波的。

图2.2.11 利用自适应中值滤波算法对加入脉冲噪声后的Lena图像处理前后对比图

图2.2.12 自适应中值滤波算法流程图

2.2.3.3 图像边缘检测算法
图像边缘是图像的基本特征之一，它包含对人类视觉和机器视觉有价值的物体图像信息。例如在利用计算机锐化图像时便有效利用了图像的边缘信息。 图像边缘检测算法是图像处理领域中一种重要的预处理技术，其广泛地应用于轮廓、特征的抽取和纹理分析等领域。(而其中较晚出现的)Canny边缘检测算法是针对局部区域或整幅图像确定一个门限[6]，对图像进行分割，从而提取出边缘，下面我们将以Canny算法为例来说明图像边缘检测算法的一种实现。

图2.2.12 应用一种边缘检测算法检测Lena图像
Canny边缘检测算法：
Canny边缘检测算法总共有四个步骤，分别是：I.采用高斯滤波的方法对原始图像作平滑处理 II.梯度幅值与方向的计算 III.对梯度幅值的非极大值抑制IV.双阈值检测和连接边缘。
I.采用高斯滤波的方法对原始图像作平滑处理：

高斯函数如上所示其中，其中 σ为高斯滤波的参数， σ的选择会直接影响滤波效果。图像平滑化的处理步骤如下：
（1）设置一个模板，在该模板内，求出所有像素灰度的加权平均值；
（2）将该值赋给该模版中心像素点的灰度值；
（3）按照同样的方法，扫描图像中的每个像素点，再进行加权平均计算。
II. 梯度幅值与方向的计算：
首先，对每个像素点求偏导数

再利用下面两个公式得到像素的梯度幅值M(i,j) 和方向 θ(i,j)。

III.对梯度幅值的非极大值抑制：
由于梯度幅值的极大值附近会产生屋脊带，故非极大值抑制的原理是通过细化幅值图中的屋脊带， 更好地确定边缘的位置非极大值抑制过程为，将梯度幅值矩阵 M(i, j) 内的每一个点都作为中心像素点，与其周围 8 个方向邻域中沿该中心点梯度方向上的两个邻域的梯度幅值进行比较。若中心像素点的幅值小于这两个方向上邻域的梯度幅值，则该中心像素点边缘标志位的值为 0。经过非极大值抑制，不仅用一个像素的宽度替代了梯度幅值矩阵的屋脊带宽度，还保留了屋脊的梯度幅值。
	经过上述三个步骤的处理后，我们已经得到了一个边缘检测的粗粒度结果，其包含有较多的伪边缘。故仍需要进一步处理。
IV.双阈值检测和连接边缘：
(1) 设定两个阈值， 即一个高阈值、一个低阈值，将梯度幅值小于低阈 值的像素点所对应的灰度值视为 0；
(2) 将低阈值 提取出的图像设为 A，将梯度幅值大于高阈值的图像 设为 B，由于 A 为通过低阈值提取得到的图像，故 边缘的连续性较好，但会含有较多的伪边缘，B 为梯度幅值大于高阈值的图像，故其不包含伪边缘， 但是其边缘的连续性不佳；
(3) 以图像 B 为基础， 图像 A 为补充，通过递归追踪法获得边缘信息。
至此，Canny算法的处理步骤便完成了。

图2.2.13利用Canny算法对两幅图片进行边缘检测的结果

2.2.3.4 传统图像处理算法存在的缺陷

	在上述的种种算法中无论是灰度变换中线性或是非线性变换中公式参数的选择还是自适应中值滤波算法中Tmax和Tmin的选择抑或是Canny边缘检测算法中双阈值的选择都具有一定的主观性。相应地，根据主观选择的不同，所导出的结果并不相同。倘若数值选择不当，则有可能会有一个糟糕的实现表现。那么，怎样去选择这些数值，怎样根据不同的实际情况去调换数值便是一个值得考虑的问题，如果一味地按照固定的参数去应用，结果可能并不尽如人意。'
	
2.2.4 传统图像处理的应用

	随着数字图像处理技术的不断发展，目前，数字图像处理的应用越来越广泛，它已经与诸多领域如遥感、公共安全、工业检测、医疗诊断、军事等深度融合，下面我们将举例说明传统图像处理的应用。
	遥感：如探索者月球图像、勇气号火星图像、再如农作物产量计算、农作物病情防治、山林防火防灾等

图2.2.4.1火星外表面图像传回
	公共安全如：
银行防盗、“天网工程”，极大保障了公民的财产和人生安全。
	工业检测如：
上文在滤波算法提到的桥梁探伤、或是道路网裂、龟裂等公路路面破损图像识别。	
	医疗诊断如：
上文在Canny边缘识别算法中提到的细胞显微成像边缘识别、又如CT技术、癌细胞识别、钙化点识别等
总结：
	传统图像处理自六十年代诞生以来，经过五十余年的不断发展，业已形成一门庞大的知识体系。其中包含了诸多算法如图像数字化、图像变换、图像增强、图像恢复、图像数据压缩、图像边缘检测、图像分割、图像特征分析、图像配准、图像融合、图像分类、图像识别、基于内容的图像检索、图形数字水印……，其广泛地应用于各种场景如地理信息科学、摄影摄像技术、航空航天、医学领域等，推动了现今社会的进步和发展。但同时，我们也应该注意到虽然传统数字图像处理可以解决诸多问题，然算法中参数的选择仍然依赖工程师的主观判断，可移植性仍有较大的进步空间。而人工智能思维的出现或许能够填补这一漏洞，使数字图像处理技术得到进一步发展。

\section{体会启发与结语}
	我负责的部分是图像的数字矩阵化。在接触此次课题之前，我对矩阵的认知是相当局限的，仅仅由线性代数的基本课程了解到它对于线性方程组与一些坐标系变换的意义，并没有体会到它如何应用于我们实际的生活中，更何况是电子计算机领域呢。通过本次课题的学习，虽然过程稍有艰辛，但我学会了不少知识，体会到了数学思维是如何作用于人类发展历史上的。在本次课题的撰写与学习过程中，我使用了C++和Python语言搭载OpenCV库来实操了一番，还是非常有趣的经历，也期待自己未来能够掌握更多的技术栈。当然，我需要更多不断地学习，需要学习更多软件工程方面的知识，才可以领会到更多的思想。
—— 谭诚

	我个人写的是传统数字图像处理的起源、发展和应用。经过这次的小组作业学习。我对数字图像处理及其发展的脉络和多样的应用场景有了大致的了解。同时也引发了个人对于数字图像处理和计算机视觉的兴趣。在和小组成员的交流沟通中，我也得到了很多锻炼和启发。虽然前后查阅资料、组织语言、制作演示文稿和其他内容的部分是不易的，但当看到小组成员汇报的结果，还是觉得颇为值得。在简要了解完传统数字图像处理之后，发觉其仍存在一定程度上的缺陷，也尝试着去利用老师上课所教授的内容去思考，是否可以用人工智能的思维去改善现有的算法，不仅拓宽了知识面，也锻炼了自己的思维，感受到大有裨益。数字图像处理是丰富多彩、颇具魅力的，其应用领域的宽广也让我切实感到技术落地的魅力，大大增强了自我不断学习相关知识、努力探索学科领域的驱动力。提升了个人对于软件工程学科的认识水平，对自我是一次全方位的发展。
—— 汪坤灿

写完神经网络与目标检测的发展史以后，最大的收获是了解了一个自己从未涉足过的领域是如何发展的。从上世纪的摸索，到最近十年间的快速发展，无数新理论、新模型、新思想在不断诞生。对于非人工智能研究者而言，人工智能的发展看起来似乎是一蹴而就的事情。但是，仅仅在目标检测这一小小的领域，就经历了一个模型从诞生到不断完善每一个部分的过程，甚至直到现在还在不断发展。因此人工智能不该被当作一个概念而热炒，似乎只要高举支持人工智能旗帜就能顺应人工智能的浪潮，而更应该了解其发展历史与具体内涵。在写神经网络与目标检测两段内容之前，我可能都不知道机器学习和深度学习有什么区别和联系，这个写作过程也相当于给自己扫了个盲。我想这份经历也可以沿用到人工智能的其他分支，乃至其他领域的学科上。除此以外，在目标检测模型的发展方面，我的内容是按照问题——解决方案——新问题的思路来撰写的，通过这种方式来理解每个看似不同的模型之间的内在联系，是一种线性连贯的解决问题方式，我想这可能也是一种比较大的启发。
—— 史晓磊

致谢
自确定开题以来，小组中的每个成员都花费了大把时间来学习新知识。面对晦涩难懂的学术论文，以及眼花缭乱的公式与知识，每个成员的努力与付出都值得被铭记，每个成员的心血与成功都值得被尊重与感谢。
同时，我们也很感谢华东师范大学软件工程学院的吴敏老师，她为我们带来了半个学期的精彩淋漓的《人工智能的数学思维》课程。很感谢她为我们介绍了近代以来的数学知识发展以及对应在电子计算机领域所诞生的新计算机科学，让我们了解到了数学思维与计算机科学的重要性。
最后，我们也非常感谢华东师范大学软件工程学院(East China Normal University, Software Engineering Institute, ECNU-SEI)开设此门课程。 

互不相容的模型组成的森林比其中每一 棵树都要睿智。

\section{ 次线性算法} 
到这里为止，我非常感兴趣的还是那些有可能获得置信度的算法的性质。我同样确定，要寻找最优预测算法，囤积海量数据是必不可少的。
然而，并不是所有数据都有价值。就像在沙漠深处安装的监视摄像头传来的视频流那样，实际上绝大部分收集来的数据毫无意义。问题在于，在这个大数据的时代，单单读取这些没有意义的数据来验证它们，真的没有意义，所需的时间就已经不切实际了——我就是这样说服自己不去查阅收到的全部邮件，这种做法没问题的!
为了解决这个问题，理论计算机学家将注意力投向了所谓的次线性算法。这些算法的特点就是在不花时间读取所有数据的情况下，仍然能够提取数据集里最核心的内容。也就是说，这些算法能够在读取输入数据时“一目十行”。
这种算法在历史上最典型的例子就是谷歌的算法。在谷歌上搜索几乎瞬间就能得到结果，然而，谷歌的数据库包含了互联网的数百万亿网页中相当大的一部分。对于谷歌来说，在向用户返回结果之前探索整个数据库根本不在考虑范围之内，因为这样做要花上几天时间!因此，谷歌的搜索算法必须是次线性的。
谷歌用到的技巧跟图书馆和词典一样，就是预先对数据库排序和整理，迅速完成对数据库的查阅。比如说，词典将单词按照字母顺序整理，就能让使用者很快知道希望查找的单词应该会在什么地方。更厉害的是，利用字母顺序，给定词典中的一页以及要查找的单词，使用者就能知道这个单词应该在词典中这一页的前面还是后面。一般来说，在(按照全序关系)排序过的数据库中查找数据非常迅速。所谓的二分算法(大概就是你用词典查单词的时候用的方法)的计算时间是数据库大小的对数。
然而谷歌、图书馆管理员和词典使用的算法都需要预先进行大量的计算。要将所有网页整理并排序，除非对其全部访问，否则不可能做到。即使如此，我们能否在既不检索整个数据集，也不预先对其进行处理的前提下，将有用的信息提取出来呢?令人惊讶的是，对于某些有用的信息以及某些数据来说，答案是肯定的。傅里叶变换就属于这种情况。
在19世纪初，约瑟夫·傅里叶等人就曾研究过傅里叶变换，它是一种对声音、图像、视频以及股票走势等信号的描述方式进行变换的方法[8]。我们在这里就只讨论声音信号。
声音可以通过鼓膜附近气压的变化来描述。但还有另一种等价的描述方式，它对于音乐的谱写来说特别简单，那就是利用组成声音的频率来描述这个声音。傅里叶变换就像一个双语词典，能够将声音的振动变化描述转换为频率描述。这样的翻译有着很多用途，一个原因是要改善声音的话，调整频率通常比调整振动更有效，另一个原因就是声音通常只由为数不多的几个频率组成。
实际上，傅里叶变换已经占据了我们的日常生活，据理查德·巴拉纽克教授所言，就连我们的计算机和电话每天都会进行数十亿次傅里叶变换。在每次听音乐、查看数字化图像，或者观看视频的时候，我们都在让机器进行傅里叶变换。
因此，计算傅里叶变换的算法如果有任何加速的可能性，都会在程序计算或者计算所需硬件的方面带来数十亿欧元的收益，更不用说用户节省下来的等待时间了。在1964年，詹姆斯·库利和约翰·图基就完成了一项壮举，大幅缩短了(离散)傅里叶变换的计算时间。尽管最简单的算法需要的时间与数据量的平方成正比，但库利和图基的算法，又叫快速傅里叶变换，可以在接近线性的时间内完成9。也就是说，要对一百万字节(1MB)的数据进行计算的话，快速傅里叶变换只需要数百万次运算(对于现代计算机来说，这相当于几毫秒的计算)，而不是朴素算法所需的100万乘以100万次运算(差不多1分钟的计算)。
9也就是说离散傅里叶变换所需的时间复杂度从$O(n^2)$变成了$O(nlogn)$。
然而在大数据的时代，快速傅里叶变换还是太慢了，尤其在处理上十亿字节(GB)甚至上万亿字节(TB)的数据时更是如此。我们能不能让傅里叶变换变得比现在更快?在2012年，哈桑尼耶、因迪克、卡塔比和普赖斯[9]发现，答案是肯定的。也可以说他们提出了一个巧妙的算法，可以对信号最主要的个频率进行近似计算，而需要的时间大概在乘以数据量的对数这个数量级10。重点在于，这个叫作稀疏傅里叶变换的精妙算法的运行时间是次线性的，也就是说它几乎忽略了整个待处理的信号，却能得知这个信号大概是什么。
10它的复杂度实际上是$O(klog^2n)$。

思考的多种模式
虽然次线性算法在计算机科学中仍不是主流，但我们的大脑似乎喜欢使用这种算法。谁又未曾一目十行读完书面材料，漫不经心地听别人讨论，或者吃饭时根本没有在意食物的味道呢?奇怪的是，有时候文章中的某些词语、讨论中的某些话题或者食物中的某些味道会吸引我们的注意力。当这种情况出现的时候，我们似乎突然就切换到了思考状态。我们会开始使用更缓慢、更精确的算法来仔细分析文章的含义、讨论的深意与本地特色菜的独特味道。
我在这里描述的，正是诺贝尔经济学奖得主丹尼尔·卡内曼的杰作《思考，快与慢》的核心。卡内曼区分了两套思考系统:系统一和系统二。系统一就像稀疏傅里叶变换，非常迅速、有效、勤奋，但有可能出现严重错误;系统二就像精确傅里叶变换，很慢且需要大量能量，懒惰但更正确。
据卡内曼所说，我们通常会将自己等同于系统二，而且几乎不会意识到系统一的存在。然而，在绝大部分时间内主导的都是系统一，而这会让我们经常犯错。试试这个:球拍和球的价格一共是1.1欧元，球拍的价格比球高1欧元，那么球的价格是多少?
很可能你立刻就看到了答案，但这个答案是错的。根据卡内曼所说，如果这种情况出现，那是因为系统一匆匆忙忙给出了第一个想到的答案;只有在之后，也许是当你看到这里的时候，你的系统二才会开始质疑系统一。
我们可能会认为卡内曼希望说服我们放弃系统一，尽可能经常让系统二运作。并非完全如此。毕竟即使系统二得出的结果更正确，但对于大脑来说，使用系统二更疲劳。思考有其代价。然而，正如所有演员、钢琴家和冲浪者都知道的那样，我们不可能一一思考自己的所有行为与动作。我们最终必须能够自然且自发地做出这些动作，无须大量有意识的思考。
要做到无须长久思考就能做出合格的动作，演员、钢琴家和冲浪者用的是同样一套方法:他们的系统二会迫使系统一学习所需的动作。对于那些希望学习或者授课的人来说，这可能是最重要的意见。与其说学习是将信息塞进脑中，不如说是让系统二来教导系统一，使得系统一发现一条捷径，能够迅速解决那些系统二已经知道怎么解决，但要花上许多时间和能量的问题。也就是说，学习的根本就是发现又快又(足够)好的捷径。

\section{迈进后严谨阶段!}
数学的学习似乎也是如此进行的。在今天，如果有人问我一个关于加法、指数或者图灵机的问题，那么对我来说，不通过系统二就能回答这些问题也并非毫无可能。这是因为，长年以来我的系统二已经成功教会系统一如何在不花费脑力劳动的前提下回答这些问题!我的大脑中储存着大量捷径，能用于解决大量数学问题，即使是那些对于年轻学生来说略感困难的问题。某些人把这种能力称为数学能力，还有些人把它称为数学直觉。我更愿意把它称为系统二通过努力发现的高效捷径。
但这不是系统一的数学训练中最重要的方面。数学家、菲尔兹奖获得者陶哲轩在他的博客中[10]将数学学习分为三个阶段，分别叫作前严谨阶段、严谨阶段和后严谨阶段。简单来说，学生首先会开始摆弄数字和数学概念，而不会忧心于自己执行的那些代数操作的有效性。然后，随着时间流逝以及接受的数学教育越来越多，严谨性的时刻就会到来，而且很快会转变为纯粹主义:一切都应该用形式化的语言做出解释。最后，后严谨阶段属于研究者，他们花上大部分时间构建不同的启发式论证，用来得出定理证明的大体轮廓，其间会暂时舍弃之前学到的形式体系。关键在于，第二个阶段似乎是到达第三个阶段的必经之路。
陶哲轩的说法可以用系统一和系统二的语言来理解。前严谨阶段对应的是系统一和系统二都没有学会严谨性的情况。当系统二发现严谨性及其重要性时，严谨阶段就开始了。在这个阶段中，系统二会教导系统一，让它意识到自身直觉的局限性。但更重要的是，系统一会由此学会测量自身的置信度，这样就能够更好地做出决定，判断自己能解决问题，还是需要向系统二求助。只有经过这种困难的学习过程，系统一才能成为系统二的完美协助者，而不会在重要的时候拖后腿。因此，成为一名优秀的数学家，基本上相当于让系统一足以在这方面胜任——但要达到这一状态，严谨阶段必不可少。
话虽如此，即使到了后严谨阶段，系统二仍会不断教导系统一，让它能够进步。这就是在莱姆基·奥利弗和孙达拉拉詹发表了关于素数最后一位数字的研究结果之后，所有数论学家的大脑中发生的事情。这个发现惊动了数论学家的系统一，它关于素数的直觉基于素数定理，而这一捷径预言相邻素数的最后一位数字大多是独立的。在这个情况
下，这个捷径出错了。自此之后，数论学家可能进行了某种近似贝叶斯推断，以减少素数定理这一捷径在相邻素数的情况中应用的置信度。
丹尼尔·卡内曼的双系统思考模型对于实用贝叶斯主义者来说的确很有趣。当然，这个模型是错的。但毕竟“所有模型都是错的”。然而卡内曼的模型似乎很有用——即使第17章暗示还存在负责创造性思考过程的第三个系统。
在实践中，尤其是在面对大数据的时候，大部分用到的算法大概是迅速的启发式算法，甚至是次线性的。然而拥有更缓慢但更正确的算法仍然是必要的。此外，如果我们拥有数个更缓慢、但我们知道足够正确的算法，那么我们就可以用这些算法来训练那些快速的启发式算法。
假如我能打个赌，(作为合格的贝叶斯主义者，我喜欢打赌!)我会说这就是未来人工智能的模样。
线性代数

- 矩阵和向量运算：在数据表示和转换中使用，如矩阵乘法、转置和逆矩阵。
- 特征值和特征向量：用于数据降维和特征选择，如主成分分析（PCA）。
- 奇异值分解（SVD）：用于数据压缩、降噪和推荐系统。
- 矩阵分解技术：如非负矩阵分解（NMF）用于主题建模。

\nocite{*}
\printbibliography[heading=bibintoc, title=\ebibname]

\newpage
